\section*{Заключение}
\addcontentsline{toc}{section}{Заключение}

В результате проведённой работы подтверждена гипотеза о том, что 
автоматизация управления очередями студентов на базе специализированного 
веб-приложения позволяет существенно повысить организованность учебного 
процесса и снизить временные затраты преподавателя на вызов студентов и 
фиксацию результатов. Теоретический анализ показал, что классическая 
FIFO-модель очереди в условиях учебной лаборатории требует адаптации под 
подтверждение присутствия студента, возможность ручного перемещения записей 
и учёт дополнительных параметров (номер лабораторной, количество работ, 
предпочтение по времени).

На основе полученных при разработке данных можно сформулировать следующие 
\textbf{выводы}:

\begin{enumerate}
    \item Использование микроархитектурного подхода на базе связки Flask 
    (Python) + Jinja2 + PostgreSQL обеспечивает необходимую гибкость, 
    позволяя реализовать ролевую модель (студент, преподаватель, 
    администратор) и CRUD-операции без усложнения клиентской логики.
    \item Проектирование модели данных с поддержкой внешних ключей и 
    каскадного удаления (PostgreSQL) является критическим условием 
    для целостности очереди, так как позволяет эффективно обрабатывать 
    записи студентов, журнал успеваемости и загруженные работы в реальном 
    времени.
    \item Внедрение WebSocket-уведомлений и механизма автоматического 
    пересчёта позиций (\texttt{normalize\_positions}) демонстрирует 
    устойчивый пользовательский эффект, выраженный в синхронном обновлении 
    очереди у всех участников и исключении конфликтов при параллельных 
    действиях.
\end{enumerate}

\textbf{Рекомендации и предложения.} Для практического применения разработанного 
решения в учебных заведениях предлагается использовать модель контейнерного 
развёртывания (Docker Compose). Это обеспечит быстрый запуск системы на 
преподавательских компьютерах или в облачной среде при минимальных затратах 
на ИТ-инфраструктуру.

\textbf{Перспективы углубления темы.} В дальнейшем исследовании планируется 
развитие функциональности системы за счёт внедрения уведомлений через 
Telegram-бота о вызове студента, модуля аналитики успеваемости (графики и 
отчёты по группам), поддержки групповых очередей (защита работы бригадой), 
а также интеграции с системами дистанционного обучения (LMS) для автоматической 
выгрузки оценок.

При разработке пояснительной записки к курсовому проекту строго соблюдались 
требования~\cite{gost_report}, определяющие структуру и правила оформления 
научно-исследовательских работ.